%%
%% Conference Paper: LS-BJOA with Hybrid CNN-LSTM
%% Fixed version — compiles cleanly on Overleaf
%%

\documentclass[sigconf]{acmart}

%% ---------------------------------------------------------------
%% Strip ALL acmart auto-generated footer/header boilerplate
%% (ACM reference format, ISBN, DOI, conference line, copyright)
%% ---------------------------------------------------------------
\setcopyright{none}
\renewcommand\footnotetextcopyrightpermission[1]{}
\pagestyle{plain}
\acmYear{2025}
%% Suppress the "ACM Reference Format" citation block on page 1
\settopmatter{printacmref=false, printccs=false, printfolios=false}

%% Remove the running header (conference name + author list strip)
\AtBeginDocument{\fancyhead{}}
\renewcommand{\headrulewidth}{0pt}

%% Short author list for page headers
\renewcommand{\shortauthors}{Gupta et al.}

%% ---------------------------------------------------------------
%% Begin document
%% ---------------------------------------------------------------
\begin{document}

%% Title
\title{LS-BJOA with Hybrid CNN-LSTM for Efficient EEG-Based Motor
Imagery Classification in Brain-Computer Interfaces}

%% ---------------------------------------------------------------
%% Authors — all share the same affiliation
%% ---------------------------------------------------------------
\author{Kunal Gupta}
\email{[KUNAL\_EMAIL]}
\affiliation{%
  \institution{Thapar Institute of Engineering and Technology}
  \city{Patiala}
  \state{Punjab}
  \country{India}
}

\author{Aman Singh}
\email{[AMAN\_EMAIL]}
\affiliation{%
  \institution{Thapar Institute of Engineering and Technology}
  \city{Patiala}
  \state{Punjab}
  \country{India}
}

\author{Karan Singh}
\email{[KARAN\_EMAIL]}
\affiliation{%
  \institution{Thapar Institute of Engineering and Technology}
  \city{Patiala}
  \state{Punjab}
  \country{India}
}

\author{Anushka Ahlawat}
\email{[ANUSHKA\_EMAIL]}
\affiliation{%
  \institution{Thapar Institute of Engineering and Technology}
  \city{Patiala}
  \state{Punjab}
  \country{India}
}

\author{Dr. Simranjit Kaur}
\email{simranjit.kaur@thapar.edu}
\affiliation{%
  \institution{Thapar Institute of Engineering and Technology}
  \city{Patiala}
  \state{Punjab}
  \country{India}
}

\author{Dr. Shubhra Dwivedi}
\email{shubhra.dwivedi@thapar.edu}
\affiliation{%
  \institution{Thapar Institute of Engineering and Technology}
  \city{Patiala}
  \state{Punjab}
  \country{India}
}

%% ---------------------------------------------------------------
%% Abstract
%% NOTE: Do NOT use underscores inside the abstract — they break
%% LaTeX. Placeholders use hyphens or braces instead.
%% ---------------------------------------------------------------
\begin{abstract}
High channel counts in EEG-based brain-computer interfaces (BCIs)
impose a significant computational burden, limiting practical
deployment of real-time assistive communication systems for patients
with amyotrophic lateral sclerosis (ALS). Existing electrode selection
techniques typically pair with classical classifiers like support
vector machines, whereas deep learning methods generally process full,
unoptimized channel configurations. No prior work has integrated
metaheuristic channel selection with deep neural networks for motor
imagery classification. We address this gap by combining the Logistic
S-shaped Binary Jaya Optimization Algorithm (LS-BJOA) with a hybrid
CNN-LSTM pipeline, evaluated on three public BCI benchmarks: BCI
Competition IV-2a (22 channels, 4-class), BCI Competition IV Dataset~1
(59 channels, binary), and BCI Competition III Dataset~IVa (118
channels, binary). The integrated system achieves accuracies of
64.29\%, 73.93\%, and 84.78\% using only
13, 33, and 62 channels,
respectively. Reducing electrode overhead while preserving
classification quality enhances feasibility for resource-constrained,
real-time BCI deployment.
\end{abstract}

%% CCS concepts --- to be added before camera-ready submission
%% \begin{CCSXML} ... \end{CCSXML}

%% Keywords
\keywords{Brain-Computer Interface, Motor Imagery, EEG Channel
Selection, LS-BJOA, CNN-LSTM, Amyotrophic Lateral Sclerosis}

\maketitle

%% ---------------------------------------------------------------
\section{Introduction}
%% ---------------------------------------------------------------

Amyotrophic lateral sclerosis (ALS) is a progressive
neurodegenerative disease that severely impairs voluntary motor
pathways, eventually leading to complete loss of muscle control.
Despite this physical deterioration, cognitive function in ALS
patients largely remains intact throughout the disease course.
Traditional assistive systems --- such as switch-based interfaces and
eye-tracking devices --- gradually become ineffective as residual
motor or ocular movement degrades in later disease stages.
Non-invasive electroencephalography (EEG) brain-computer interfaces
(BCIs) provide a viable alternative by capturing scalp electrical
activity directly, enabling communication and environmental control
without any muscular exertion~\cite{wolpaw2002bci}.

A persistent challenge in EEG-based motor imagery (MI) systems is
the high dimensionality of the recorded signal, which routinely spans
dozens to over a hundred electrodes. Many of these channels carry
redundant or task-irrelevant information, increasing computational
overhead, slowing inference, and often degrading classifier
performance. Metaheuristic channel selection algorithms reduce
electrode counts effectively, but they are typically paired with
classical classifiers such as support vector machines (SVM) or linear
discriminant analysis (LDA)~\cite{tiwari2023lsbjoa}. These models
struggle with the highly non-linear spatiotemporal structure of EEG
signals. Meanwhile, modern deep learning architectures such as
EEGNet~\cite{lawhern2018eegnet} process the entire electrode layout
without any channel optimization step, leaving both problems
unresolved independently.

To close this gap, we integrate the Logistic S-shaped Binary Jaya
Optimization Algorithm (LS-BJOA), as proposed by Tiwari et
al.~\cite{tiwari2023lsbjoa}, with a hybrid deep learning classifier
combining a convolutional neural network (CNN) and a long short-term
memory (LSTM) network. LS-BJOA selects the most informative electrode
subset; the CNN extracts local spatial patterns from the active
channels; and the LSTM captures sequential temporal dynamics across
the trial. We evaluate this pipeline on the three public benchmark
datasets used in the original LS-BJOA study, enabling direct
comparison against its SVM baseline. We also benchmark six deep
learning architectures --- EEGNet, CNN, RNN, LSTM, CNN-RNN, and
CNN-LSTM --- under both full-channel and channel-selected conditions.

The rest of this paper is organized as follows:
Section~\ref{sec:datasets} describes the benchmark datasets.
Section~\ref{sec:method} presents the proposed methodology.
Section~\ref{sec:results} reports experimental results and discussion.
Section~\ref{sec:conclusion} concludes the paper.

%% ---------------------------------------------------------------
\section{Datasets}
\label{sec:datasets}
%% ---------------------------------------------------------------

We evaluate our proposed pipeline on three publicly available BCI
benchmark datasets. Using the same datasets as the baseline
study~\cite{tiwari2023lsbjoa} allows direct comparison against the
original SVM-based classifier. Together, these datasets span a range
of electrode configurations, task complexities, and trial structures,
covering both binary and four-class motor imagery paradigms.

\textbf{BCI Competition IV Dataset~2a}~\cite{brunner2008bci4} contains
motor imagery recordings from 9 healthy subjects. The setup uses 22
EEG channels and 3 EOG channels, sampled at 250~Hz. Subjects
performed four motor imagery tasks: left hand, right hand, feet, and
tongue. Each subject completed two sessions of 6 runs each, yielding
288 trials per session. Raw signals were stored in GDF format and
bandpass filtered between 0.5 and 100~Hz with a notch at 50~Hz.

\textbf{BCI Competition IV Dataset~1}~\cite{tangermann2012bci4}
comprises recordings from 7 healthy participants across 59 EEG
channels sampled at 100~Hz. The experimental design is binary: some
participants performed left hand and foot imagery while others
performed right hand and left hand imagery. The dataset contains two
calibration runs of 100 trials each; the first run is split 80/20 for
training and testing.

\textbf{BCI Competition III Dataset~IVa}~\cite{blankertz2006bci3}
includes data from 5 subjects (aa, al, av, aw, ay) recorded over
118 EEG channels at 100~Hz. The binary task involves distinguishing
left hand from right hand motor imagery. Each trial presented a visual
cue for 3.5~s. Each subject performed 140 trials in total, with 110
allocated for training and 30 for validation. The large electrode
count makes this dataset the primary testbed for our channel reduction
algorithm.

Table~\ref{tab:datasets} summarizes the key characteristics of each
benchmark dataset.

\begin{table}[t]
  \caption{Summary of benchmark EEG datasets used for evaluation.}
  \label{tab:datasets}
  \begin{tabular}{lrrrcl}
    \toprule
    Dataset      & Subj. & Ch. & Classes & Hz  & Task       \\
    \midrule
    BCI IV-2a    & 9     & 22  & 4       & 250 & Multi-class \\
    BCI IV-1     & 7     & 59  & 2       & 100 & Binary      \\
    BCI III-IVa  & 5     & 118 & 2       & 100 & Binary      \\
    \bottomrule
  \end{tabular}
\end{table}

%% ---------------------------------------------------------------
\section{Methodology}
\label{sec:method}
%% ---------------------------------------------------------------
We propose a complete closed-loop brain-computer interface system that utilizes metaheuristic channel selection coupled with a hybrid deep learning classifier. The conceptual data flow and system architecture are illustrated in Figure~\ref{fig:architecture}. The pipeline starts with raw multichannel EEG signals, preprocesses them to remove noise and eye-movement artifacts, optimizes the electrode configuration using LS-BJOA, feeds the masked channels into a hybrid CNN-LSTM network, and outputs a classification decision to control assistive interfaces or communication tools.

\begin{table*}[t]
  \caption{Classification results on BCI Competition IV Dataset 2a (4-Class, 22 channels).}
  \label{tab:results_2a}
  \centering
  \begin{tabular}{lccccc}
    \toprule
    & \multicolumn{2}{c}{Full Channels (22)} & \multicolumn{3}{c}{Selected Channels} \\
    \cmidrule(r){2-3} \cmidrule(l){4-6}
    Model & Acc (\%) & F1-Score & Acc (\%) & F1-Score & Active Ch. \\
    \midrule
    EEGNet & 45.40 & 0.443 & 45.02 & 0.445 & 13 \\
    CNN & 40.80 & 0.400 & 38.31 & 0.372 & 13 \\
    RNN & 48.89 & 0.544 & 58.51 & 0.572 & 13 \\
    LSTM & 51.80 & 0.497 & 50.84 & 0.587 & 13 \\
    CNN+RNN & 33.14 & 0.232 & 27.39 & 0.171 & 13 \\
    CNN+LSTM & 55.63 & 0.515 & 64.29 & 0.577 & 13 \\
    \bottomrule
  \end{tabular}
\end{table*}
\begin{table*}[t]
  \caption{Classification results on BCI Competition IV Dataset 1 (Binary, 59 channels).}
  \label{tab:results_ds1}
  \centering
  \begin{tabular}{lccccccc}
    \toprule
    & \multicolumn{3}{c}{Full Channels (59)} & \multicolumn{4}{c}{Selected Channels} \\
    \cmidrule(r){2-4} \cmidrule(l){5-8}
    Model & Acc (\%) & F1-Score & ROC-AUC & Acc (\%) & F1-Score & ROC-AUC & Active Ch. \\
    \midrule
    EEGNet & 74.76 & 0.723 & 0.800 & 73.93 & 0.744 & 0.780 & 33 \\
    CNN & 67.98 & 0.685 & 0.753 & 70.60 & 0.706 & 0.759 & 33 \\
    RNN & 67.14 & 0.681 & 0.745 & 66.67 & 0.642 & 0.717 & 33 \\
    LSTM & 64.52 & 0.632 & 0.690 & 61.19 & 0.518 & 0.657 & 33 \\
    CNN+RNN & 67.26 & 0.715 & 0.747 & 70.71 & 0.713 & 0.769 & 33 \\
    CNN+LSTM & 66.07 & 0.603 & 0.716 & 66.79 & 0.624 & 0.728 & 33 \\
    \bottomrule
  \end{tabular}
\end{table*}
\begin{table*}[t]
  \caption{Classification results on BCI Competition III Dataset IVa (Binary, 118 channels).}
  \label{tab:results_ds3}
  \centering
  \begin{tabular}{lccccccc}
    \toprule
    & \multicolumn{3}{c}{Full Channels (118)} & \multicolumn{4}{c}{Selected Channels} \\
    \cmidrule(r){2-4} \cmidrule(l){5-8}
    Model & Acc (\%) & F1-Score & ROC-AUC & Acc (\%) & F1-Score & ROC-AUC & Active Ch. \\
    \midrule
    EEGNet & 82.84 & 0.779 & 0.907 & 84.78 & 0.785 & 0.896 & 62 \\
    CNN & 78.85 & 0.747 & 0.859 & 83.26 & 0.799 & 0.899 & 62 \\
    RNN & 79.39 & 0.713 & 0.858 & 80.34 & 0.748 & 0.892 & 62 \\
    LSTM & 60.41 & 0.554 & 0.617 & 62.59 & 0.585 & 0.640 & 62 \\
    CNN+RNN & 79.98 & 0.763 & 0.853 & 79.63 & 0.724 & 0.918 & 62 \\
    CNN+LSTM & 71.58 & 0.632 & 0.771 & 74.55 & 0.575 & 0.709 & 62 \\
    \bottomrule
  \end{tabular}
\end{table*}

\begin{figure}[htbp]
  \centering
  \includegraphics[width=0.85\linewidth]{architecture}
  \caption{Architecture of the proposed BCI system incorporating LS-BJOA channel selection and hybrid deep learning classification.}
  \label{fig:architecture}
\end{figure}

\subsection{Preprocessing Pipeline}
To ensure the model learns robust neurological patterns rather than ambient noise or ocular artifacts, the raw EEG recordings are passed through a sequential preprocessing pipeline:
\begin{enumerate}
    \item \textbf{Filtering:} A high-pass filter at 1~Hz is applied to eliminate slow DC drift, followed by a bandpass filter (typically 4--40~Hz) to isolate the motor imagery mu and beta rhythms.
    \item \textbf{Artifact Rejection:} Independent Component Analysis (ICA) is computed on the filtered signals. Independent components corresponding to eye blink and movement artifacts are automatically identified using electrooculography (EOG) reference channels and zeroed out before reconstructing the signal.
    \item \textbf{Trial Segmentation:} The continuous signals are segmented into individual motor imagery epochs ($T_{\text{samples}}$).
    \item \textbf{Per-Trial Normalization:} To prevent data leakage and cross-trial bias, Z-score normalization is applied independently to each individual trial. For a given trial matrix $X \in \mathbb{R}^{C \times T}$, the normalized value of channel $c$ at time sample $t$ is computed as:
    \begin{equation}
        X_{\text{norm}}(c, t) = \frac{X(c, t) - \mu_c}{\sigma_c}
    \end{equation}
    where $\mu_c$ and $\sigma_c$ are the mean and standard deviation calculated exclusively over the time samples $T$ of channel $c$ within that single trial.
\end{enumerate}

\subsection{Channel Selection}
The S-shaped Binary Jaya Optimization Algorithm (LS-BJOA) is utilized to search for the optimal subset of electrodes. Unlike other metaheuristics, Jaya does not require algorithm-specific control parameters. In standard Jaya, the candidate solution vector is updated by moving toward the best solution and away from the worst:
\begin{equation}
    X'_{i, j} = X_{i, j} + r_1 \left(X_{\text{best}, j} - |X_{i, j}|\right) - r_2 \left(X_{\text{worst}, j} - |X_{i, j}|\right)
\end{equation}
where $X_{i, j}$ represents the $j$-th decision variable of candidate $i$, $X_{\text{best}, j}$ is the best solution in the population, $X_{\text{worst}, j}$ is the worst, and $r_1, r_2 \in [0, 1]$ are uniform random numbers.

To adapt the algorithm to the discrete space of channel selection, the continuous updates are mapped to binary decisions $\{0, 1\}$ using a logistic S-shaped transfer function:
\begin{equation}
    T\left(X'_{i, j}\right) = \frac{1}{1 + e^{-X'_{i, j}}}
\end{equation}
The binary state of the electrode mask is then updated as:
\begin{equation}
    X^{\text{binary}}_{i, j} = \begin{cases} 
      1, & \text{if } r_3 < T\left(X'_{i, j}\right) \\
      0, & \text{otherwise}
   \end{cases}
\end{equation}
where $r_3 \in [0, 1]$ is a random number.

\paragraph{Fitness Evaluation:}
To evaluate the quality of a selected electrode subset during optimization, a fast proxy classifier is trained to avoid expensive deep learning training cycles. A linear Support Vector Classifier (LinearSVC) is trained on simple log-variance features extracted from the selected channels. The objective is to maximize classification accuracy while minimizing the number of active channels. The fitness function is formulated as:
\begin{equation}
    \text{Fitness} = \alpha \cdot \text{Accuracy} + (1 - \alpha) \cdot \left(1 - \frac{N_{\text{selected}}}{N_{\text{total}}}\right)
\end{equation}
where $N_{\text{selected}}$ is the number of active channels, $N_{\text{total}}$ is the total available channels, and $\alpha = 0.9$ is a weighting coefficient that prioritizes classification performance.

\subsection{CNN-LSTM Classifier}
Once the optimal binary channel mask is determined, only the selected EEG channels are fed into the deep learning pipeline. The hybrid network comprises:
\begin{itemize}
    \item \textbf{Spatial Convolutional Layers:} 2D convolutional layers with kernels that perform spatial filtering across the selected electrodes to capture localized activation patterns.
    \item \textbf{Temporal Recurrent Layers:} LSTM layers that process the sequential spatial feature maps over time. The recurrent units capture long-term temporal dependencies essential for modeling the motor imagery planning phase.
    \item \textbf{Classification Head:} Fully connected layers with dropout regularization, culminating in a softmax activation function to output class probabilities.
\end{itemize}

\subsection{Experimental Setup}
The deep learning models are optimized using the Adam optimizer with a cross-entropy loss function. The initial learning rate is set to $10^{-3}$ and is dynamically reduced when the validation loss plateaus. Early stopping with a patience of 5 epochs is implemented to prevent overfitting. Models are trained using $K$-fold cross-validation on the preprocessed training split and evaluated on the held-out test split.

%% ---------------------------------------------------------------
\section{Results and Discussion}
\label{sec:results}
%% ---------------------------------------------------------------
We benchmarked six deep learning architectures under two conditions: (1) using the full electrode configuration, and (2) using the optimized channel subset selected by LS-BJOA. The evaluation metrics include classification Accuracy and F1-score across all datasets, as well as area under the ROC curve (ROC-AUC) for the binary paradigms.

Tables~\ref{tab:results_2a}, \ref{tab:results_ds1}, and \ref{tab:results_ds3} present the experimental results for the three benchmark datasets.

For the 4-class BCI Competition IV Dataset 2a, the classification performance under the full 22-channel configuration varies across architectures. The hybrid CNN-LSTM model achieves the highest accuracy of 55.63\% with an F1-score of 0.515. When using the LS-BJOA optimized channel configuration, which reduces the active channel count to 13 (a 41\% reduction in electrodes), the CNN-LSTM model's test accuracy increases to 64.29\% and its F1-score rises to 0.577. Similarly, the RNN model benefits from the channel selection step, increasing in accuracy from 48.89\% (full) to 58.51\% (selected). This performance enhancement indicates that discarding redundant or noisy electrodes mitigates spatial overfitting, acting as a spatial regularizer that assists recurrent structures in capturing sequence dynamics.

A similar trend is observed on the binary datasets with larger initial electrode configurations. On the 59-channel BCI Competition IV Dataset 1, the EEGNet model achieves the highest accuracy of 74.76\% under the full configuration. After LS-BJOA selection, which discards 26 channels (reducing the active set to 33 channels), the EEGNet classifier achieves a comparable accuracy of 73.93\% while maintaining similar F1-score and ROC-AUC profiles. On the 118-channel BCI Competition III Dataset IVa, the effect of channel selection is more pronounced. The full-channel EEGNet yields 82.84\% accuracy and 0.907 ROC-AUC, which improves to 84.78\% accuracy using only 62 channels. Similarly, the CNN accuracy increases from 78.85\% to 83.26\% with the same 62-channel mask. This indicates that high-density configurations contain considerable spatial redundancy that, when unoptimized, introduces non-task-related noise into the convolutional layers.

Across all three datasets, the hybrid CNN-LSTM and the compact EEGNet architectures consistently outperform simple RNN and LSTM networks. In the full configuration, the spatial convolutions in CNN-LSTM and EEGNet capture structural patterns across the electrode grid before temporal modeling, whereas pure recurrent networks process the raw signal sequences directly and struggle to extract localized features. For example, on Dataset IVa, the simple LSTM model achieves only 60.41\% accuracy, compared to the 82.84\% achieved by EEGNet. While channel selection improves simple recurrent network accuracy (such as the LSTM model on Dataset IVa rising to 62.59\%), the hybrid CNN-LSTM and EEGNet networks maintain their performance lead, demonstrating the utility of combining localized spatial filtering with temporal sequence learning.

\section{Conclusion}
\label{sec:conclusion}
We integrated the Logistic S-shaped Binary Jaya Optimization Algorithm (LS-BJOA) with a hybrid CNN-LSTM network to optimize electrode configurations for motor imagery classification. The framework was evaluated on three benchmark datasets: BCI Competition IV-2a, BCI Competition IV Dataset 1, and BCI Competition III Dataset IVa. The experimental results demonstrate that metaheuristic electrode reduction maintains or improves classification accuracy while significantly reducing the number of active channels. Specifically, the hybrid CNN-LSTM model achieved a test accuracy of 64.29\% on BCI Competition IV-2a using only 13 of the 22 channels, outperforming the full-channel configuration by 8.66\%. Similarly, the EEGNet model achieved 84.78\% accuracy on BCI Competition III Dataset IVa using 62 of the 118 channels. These findings indicate that removing redundant electrodes filters out spatial noise, which helps the classifier avoid overfitting. Future work will investigate cross-subject generalization to evaluate the robustness of the selected electrode subsets across different participants.

%% ---------------------------------------------------------------
%% Bibliography
%% ---------------------------------------------------------------
\bibliographystyle{ACM-Reference-Format}
\bibliography{references}

\end{document}
